\begin{table}[htbp]
\centering
\caption{RAs com mais de um segmento: grupo de cada segmento}
\label{tab:ra_segmentos}
\small
\begin{tabular}{l c c p{5.6cm}}
\hline
\textbf{RA} & \textbf{Segmentos} & \textbf{Grupos} & \textbf{Grupo de cada segmento} \\ \hline
Paranoá & 4 & 2 & AE (B); AF (B); AS (A); ML (B) \\
Planaltina & 4 & 2 & AE (B); AF (C); AS (B); ML (B) \\
Brazlândia & 3 & 2 & AF (C); AS (C); ML (A) \\
Gama & 3 & 2 & AE (C); AF (C); ML (A) \\
Samambaia & 3 & 2 & AE (C); AF (B); ML (B) \\
Sobradinho II & 3 & 2 & AE (C); AF (C); ML (A) \\
Park Way & 2 & 2 & AF (A); AS (C) \\
Recanto das Emas & 2 & 2 & AE (C); AF (B) \\
Riacho Fundo II & 2 & 2 & AS (C); ML (A) \\
Sobradinho & 2 & 2 & AE (C); ML (B) \\
Taguatinga & 2 & 2 & AE (C); ML (A) \\
\hline
Fercal & 3 & 1 & AE (B); AF (B); AS (B) \\
São Sebastião & 3 & 1 & AE (B); AS (B); ML (B) \\
Ceilândia & 2 & 1 & AE (C); AF (C) \\
Riacho Fundo & 2 & 1 & AF (A); ML (A) \\
Sol Nascente/Pôr do Sol & 2 & 1 & AS (B); ML (B) \\
\hline
\end{tabular}
\par\footnotesize Fonte: elaboração própria a partir da PDAD Rural 2022 (IPEDF). AF: Agricultura Familiar; AE: Agricultura Empresarial; AS: Assentamento/Agrovila; ML: Moradia e Lazer. A a C: grupos do $K$-Means ($K = 3$). As 7 RAs com um único segmento não aparecem.
\end{table}